\documentclass{../tex/report}
\usepackage{setspace} % Setting line spacing
\usepackage{ulem} % Underline
\usepackage{caption} % Captioning figures
\usepackage{subcaption} % Subfigures
\usepackage{geometry} % Page layout
\usepackage{multicol} % Columned pages
\usepackage{array,etoolbox}
\usepackage{fancyhdr}
\usepackage{enumitem}
\usepackage[toc,page]{appendix}
\setlist{noitemsep}

% Page layout (margins, size, line spacing)
\geometry{letterpaper, left=1in, right=1in, bottom=1in, top=1in}
\setstretch{1}

% Headers
\pagestyle{fancy}
\lhead{I${}^2$CIP API}
\rhead{PeaPod Technologies Inc.}

\begin{document}

\begin{titlepage}
    \begin{center}
        \vspace*{1.2cm}

        \textbf{\large{I${}^2$CIP API}}

        \vspace{0.5cm}

        A Hardware Design Specification for Bus-Switched Hot-Swap Modules of I${}^2$C Targets, and a Software Library of Intra-Network Communications Protocols for Rapid Implementation of Plug-and-Play Embedded Systems

        \vfill
        \scriptsize{
    \textbf{Jayden Lefebvre - Founder \& CEO, Lead Engineer}\\
    Northumberland County, ON, Canada\\
    \vspace{.5cm}
    \textbf{Nathan Chareunsouk - Co-Founder \& CFO, Design Lead}\\Greater Toronto Area, ON, Canada\\
    \vspace{.5cm}
    \textbf{Grayson Closs - Software Engineer}
}

\vspace{.75cm}

Primary Contact Email: \uline{contact@peapodtech.com}
        \vspace{1.25cm}

        Revision 0.1\\
        PeaPod Technologies Inc.\\
        October 27th, 2026

    \end{center}
\end{titlepage}

\thispagestyle{plain}

\tableofcontents
\clearpage

\section{Product Overview}

% ## Switching Multiplexers

% On an I2C bus, only one device may occupy any given 7-bit I2C address at a time. By splitting the I2C bus into many individually-selectable sub-busses/"subnetworks" using a multiplexer ("MUX"), one is able to communicate with devices that share I2C address space independently on multiple subnetworks.

% MUX bus 0 is reserved as the "default bus"; SPRT EEPROM **must** be located on this bus.

% MUX bus 7 is reserved as the "fake bus"; no devices may be located on this bus.

% ## Fully Qualified Addressing (FQA)

% > **FQA** defines 3 new prefixed address segments for addressing devices across all available I2C intra-networks; a necessary expansion on the I2C *device-address* system.

% 1. The most significant 3 bits are reserved to indicate which of up to 8 I2C busses this intra-network is on.
% 2. There are 8 reserved multiplexer addresses: `0x70` → `0x77`; the following 3 bits are used to encode the "module number", 0 through 7, indicating which multiplexer is being addressed.
% 3. There are 8 busses available on each multiplexer; 3 bits are used to encode the bus.
% 4. There are 7 bits of uniquely-addressable *device-address* space on each I2C bus.

% **FQA size: 16 bits.** Together, segments 1 through 3 (bits 7 through 15) define the **subnetwork**, while segment 4 (bits 0 through 6) still define the **device-address**.

% Segment Extraction: `(address >> LSB) & (0xFFFF >> (16 - LEN))`

% Visualized in binary: `0bNETMODBUSADDRESS`

% In decimal, like an IP address: `N : M : B : ADR`.

% E.g. `0:3:1:043` (`0b000:011:001:0101011`, `0x0668`) indicates a device on I2C bus 0, module 3 (MUX address 0x73), MUX bus 1 (MUX instruction `0b00000010`), device address 43 (`0x2B`).

% ## Protocols

% | I2CIP Protocol         | IP Equivalent | Function                                                                                                         | Library Implementation                                                                                    | OSI Layer(s)       |
% | ---------------------- | ------------- | ---------------------------------------------------------------------------------------------------------------- | --------------------------------------------------------------------------------------------------------- | ------------------ |
% | Device Lookup Protocol | ARP           | Allows device groups to be looked up by ID. Informs "keep-alive" in a rebuild.																		| Routing Table (Hash Table): Device ID → FQA [ ]                                           								| Data Link, Network |
% | Reverse DL Protocol    | RARP          | Allows device IDs to be looked up by FQA.																																				| Routing Table (BST): FQA → Device ID 																																			| Data Link, Network |
% | Network Scan Protocol  | RIP           | Builds the Routing Table. Uses SPRT EEPROM and any local SPRT to attempt address recognition (otherwise ignore). | Ping for modules + One-time local all-scan → Read EEPROM & One-time module all-scan → Build Routing Table | Data Link, Network |
% | Bus Switching Protocol | SNMP          | MUX bus switching                                                                                                | Write `1 << BUS` to MUX                                                                                   | Network            |
% |                        |               |                                                                                                                  |                                                                                                           |                    |

% > Note: ARP and RARP are two different traversal methods of the same tree. This is done in order to minimize heap memory allocation (necessary on microcontroller architectures).

% # Compatible Hardware Specification

% ## Multiplexers

% > Multiplexers are used to split the I2C bus into many individually-selectable sub-busses, or simply busses. This allows for multiple devices to share I2C/device-address space independently on multiple subnets, enabling massive scalability of dynamic routing.

% The I2C Multiplexer is a device that allows for the selection of one of 8 busses to be connected to the I2C bus. It is selectable by a 3-bit address, 0x70 to 0x77, and has 8 busses. Busses are selected in a bitwise fashion: `0b00000001` selects bus 0, `0b00000010` selects bus 1, etc. while disabling all other busses, and `0b00000000` disables all busses.

% ## Static Partial Routing Table (SPRT) EEPROM

% > 24LC32 16-bit 4kB EEPROM on Bus 0 at address `0x50` stores a UTF-8-encoded JSON-formatted partial routing table (with *decimal-encoded* addresses).

% A partial routing table (see below) indexes only the devices on a bus.

% ```json
% // Indexed by BUS
% [
% 	// Bus 0 (default bus)
% 	{
% 		// Array is used to list all occurrences of flexible-address devices
% 		"eeprom" : [ 80 ], // Decimal encoding (0x50) - This is the EEPROM we are reading from
% 		"mydevice" : [ 101, 43 ],
% 	},
% 	//Bus 1, etc. up to Bus 7 (makes 8 buses)
% 	{ },
% 	{
% 		"eeprom" : [80], // A different eeprom, on bus 2
% 	},
% 	{ },
% 	{ },
% 	{ },
% 	{ },
% 	{ },
% ]
% ```

\clearpage

% References
\bibliographystyle{IEEEtran}
\bibliography{references}
\end{document}